\documentclass[11pt,a4paper]{article}

\usepackage[utf8]{inputenc}
\usepackage[T1]{fontenc}
\usepackage{amsmath,amssymb,amsthm}
\usepackage{graphicx}
\usepackage{booktabs}
\usepackage{multirow}
\usepackage{hyperref}
\usepackage{algorithm}
\usepackage{algorithmic}
\usepackage{xcolor}
\usepackage{subcaption}
\usepackage[margin=1in]{geometry}
\usepackage{listings}
\usepackage{natbib}
\usepackage{enumitem}

\newtheorem{definition}{Definition}
\newtheorem{proposition}{Proposition}
\newtheorem{remark}{Remark}

\lstset{
    basicstyle=\ttfamily\small,
    breaklines=true,
    frame=single,
    language=Python,
    numbers=left,
    numberstyle=\tiny,
}

\title{Deductive-State MCTS with Symbolic Verification\\for Mathematical Reasoning}
\author{
    Author Name\\
    Affiliation\\
    \texttt{email@example.com}
}
\date{}

\begin{document}
\maketitle

% ===========================================================================
\begin{abstract}
% ===========================================================================

Large language models for mathematical reasoning frequently produce solutions
containing symbolic contradictions and code execution errors that go undetected
during standard supervised fine-tuning. We propose Deductive-State Monte Carlo
Tree Search (DS-MCTS), a training and inference framework that integrates
symbolic verification via SymPy and code execution checking directly into the
MCTS expansion phase. Unlike standard MCTS approaches that rely on learned value
functions or outcome-based rewards, DS-MCTS uses deterministic symbolic and
runtime verification to prune logically inconsistent reasoning branches before
they are expanded, reducing wasted computation and focusing the search on
deductively sound solution paths. We introduce a two-phase training pipeline:
supervised fine-tuning on the OpenMathReasoning dataset establishes a strong
baseline, followed by MCTS-guided rejection sampling that fine-tunes the model
exclusively on solutions that survive both symbolic and code verification. We
evaluate DS-MCTS across three 7B-parameter base models---NuminaMath-7B-TIR,
OpenMath-Nemotron-7B, and Qwen2.5-Math-7B-Instruct---demonstrating that the
method generalizes across architectures and pre-training regimes. Comprehensive
ablations over verification components, MCTS hyperparameters, problem
categories, and difficulty tiers reveal that composite verification yields the
largest accuracy gains on olympiad-level problems, where models produce the
highest rate of symbolic contradictions. All code, trained adapters, and
experimental results are publicly available.

\end{abstract}


% ===========================================================================
\section{Introduction}
\label{sec:introduction}
% ===========================================================================

Recent advances in large language models (LLMs) have produced remarkable
progress on mathematical reasoning benchmarks. Models such as GPT-4
\citep{openai2023gpt4}, Gemini \citep{google2024gemini}, and open-weight
alternatives including Qwen2.5-Math \citep{yang2024qwen25math} and
DeepSeek-Math \citep{shao2024deepseekmath} have demonstrated the ability to
solve competition-level mathematics problems that were considered beyond the
reach of automated systems only a few years ago. The AI Mathematical Olympiad
(AIMO) Progress Prize competitions have served as a focal point for this
progress: the NuminaMath team won the first AIMO Progress Prize in 2024 using
tool-integrated reasoning (TIR) with a 7B-parameter model
\citep{numina2024winning}, and NVIDIA secured first place in AIMO-2 using
models trained on the OpenMathReasoning dataset \citep{moshkov2025aimo2}.

Despite these advances, LLMs for mathematical reasoning exhibit a fundamental
weakness: they produce solutions that \emph{appear} correct---with well-formed
reasoning chains, syntactically valid code, and plausible final answers---but
contain internal contradictions that a symbolic checker could detect. A model
might establish that $x = 5$ in one reasoning step and proceed to compute
$2x = 12$ in the next, or produce a Python code block that raises a
\texttt{NameError} when executed. Standard supervised fine-tuning (SFT) treats
all training solutions equally, including those with such latent errors. The
model learns to reproduce the surface patterns of correct solutions without
acquiring the ability to detect or avoid symbolic inconsistency.

Existing approaches to improving LLM reasoning quality have significant
limitations in this regard. Chain-of-thought (CoT) prompting
\citep{wei2022chain} encourages step-by-step reasoning but provides no mechanism
for verifying that the steps are mutually consistent. Self-consistency
\citep{wang2023selfconsistency} generates multiple independent solutions and
selects the majority answer, but fails when the model systematically produces
the same type of error across samples. Tool-integrated reasoning (TIR)
\citep{numina2024winning} interleaves code execution within the reasoning
chain, which catches runtime errors in individual code blocks, but does not
verify symbolic consistency across reasoning steps or between natural language
claims and code outputs. Process reward models (PRMs)
\citep{lightman2023lets} learn to score individual reasoning steps, but require
expensive human annotation and provide probabilistic rather than deterministic
verification signals.

We propose \textbf{Deductive-State Monte Carlo Tree Search (DS-MCTS)}, a
framework that addresses these limitations by integrating deterministic
symbolic and runtime verification directly into the tree search process. The
key insight is that mathematical reasoning steps make claims that can be
checked against each other: equations can be parsed and evaluated for
consistency using computer algebra systems, and code blocks can be executed to
verify they run without errors. By performing these checks at each node
expansion in the MCTS tree, we prune deductively invalid branches before the
search wastes computation exploring them.

DS-MCTS operates at two levels. At \emph{inference time}, it guides the search
toward symbolically consistent solution paths, improving accuracy without
retraining. At \emph{training time}, it implements a form of rejection sampling:
the model generates candidate solutions via tree search, only those surviving
verification are retained, and the model is fine-tuned on this curated set. This
MCTS-guided rejection sampling (which we term MCTS-RL) produces a model that
has internalized the verification signal---it generates fewer symbolic
contradictions and code errors even when used without MCTS at inference time.

Our contributions are:

\begin{enumerate}[nosep]
    \item \textbf{DS-MCTS algorithm}: We formalize MCTS with integrated symbolic
    (SymPy) and code execution verification during tree expansion, defining the
    notion of a \emph{deductive state} as the accumulated, verifiable reasoning
    context at each tree node.

    \item \textbf{Two-phase training pipeline}: We demonstrate that SFT followed
    by MCTS-guided rejection sampling (MCTS-RL) consistently outperforms SFT
    alone across multiple base models.

    \item \textbf{Multi-model generalizability}: We evaluate across three
    7B-parameter models with different pre-training regimes---NuminaMath-7B-TIR,
    OpenMath-Nemotron-7B, and Qwen2.5-Math-7B-Instruct---showing that DS-MCTS
    is not architecture-specific.

    \item \textbf{Comprehensive ablation}: We provide detailed analysis of
    verification components (symbolic vs. code vs. composite), MCTS
    hyperparameters (depth, branching, UCT constant, iterations), problem
    categories, and difficulty tiers, identifying when and why verification
    helps most.
\end{enumerate}

The remainder of this paper is organized as follows. Section~\ref{sec:related}
reviews related work on mathematical reasoning, MCTS for LLMs, and symbolic
verification. Section~\ref{sec:method} presents the DS-MCTS algorithm, the
verification components, and the two-phase training pipeline in full detail.
Section~\ref{sec:setup} describes the experimental setup including models,
dataset, and evaluation methodology. Section~\ref{sec:results} presents the
main results and ablation studies. Section~\ref{sec:analysis} provides deeper
analysis of tree structure, failure modes, and token efficiency.
Section~\ref{sec:discussion} discusses implications and limitations, and
Section~\ref{sec:conclusion} concludes.


% ===========================================================================
\section{Related Work}
\label{sec:related}
% ===========================================================================

\subsection{Mathematical Reasoning with Large Language Models}

The application of LLMs to mathematical reasoning has progressed rapidly along
several axes. Chain-of-thought prompting \citep{wei2022chain} demonstrated that
eliciting intermediate reasoning steps substantially improves performance on
arithmetic and word problems. Self-consistency decoding
\citep{wang2023selfconsistency} extended this by sampling multiple CoT paths and
selecting the answer via majority voting, exploiting the observation that
correct reasoning paths tend to converge on the same answer. However,
self-consistency is fundamentally limited by the assumption that errors are
random and independent; when a model systematically produces a particular type
of symbolic error, majority voting reinforces rather than corrects it.

Tool-integrated reasoning (TIR) represents a different approach: instead of
relying solely on the model's internal computation, TIR-format solutions
interleave natural language reasoning with executable Python code blocks. The
code is executed by an external interpreter, and the output is fed back into the
reasoning chain. NuminaMath \citep{numina2024winning} demonstrated the
effectiveness of this approach, winning the first AIMO Progress Prize using a
7B-parameter model fine-tuned on TIR-format solutions. The OpenMathReasoning
dataset \citep{nvidia2025openmath} subsequently provided 1.7 million TIR-format
solutions generated by frontier models, enabling large-scale SFT on this format.

Despite TIR's integration of code execution, a gap remains: TIR verifies that
individual code blocks execute without runtime errors, but does not check
symbolic consistency \emph{across} reasoning steps. A model can correctly
execute a code block that computes $x = 5$ and then, in the next natural
language reasoning step, assert that $x = 3$ without any mechanism detecting
the contradiction. Our work addresses this gap by adding cross-step symbolic
verification alongside code execution checking.

\subsection{Monte Carlo Tree Search for Language Models}

MCTS has a long history in game-playing AI, from its origins in Go
\citep{coulom2006efficient, kocsis2006bandit} to its integration with deep
neural networks in AlphaGo and AlphaZero \citep{silver2017mastering}. Its
application to language model reasoning is more recent. Tree-of-Thoughts
\citep{yao2023tree} introduced deliberate tree-structured exploration of
reasoning paths, using the LLM itself as a value function to evaluate partial
solutions. However, Tree-of-Thoughts relies on the model's own judgment of
solution quality, which is precisely the judgment that produces errors in the
first place.

AlphaProof \citep{deepmind2024alphaproof} applied MCTS to formal mathematical
theorem proving, achieving silver-medal performance at the International
Mathematical Olympiad. AlphaProof operates in the formal proof domain (Lean),
where verification is complete and sound: a proof either type-checks or it does
not. Our work occupies a middle ground between the heavyweight formal
verification of AlphaProof and the no-verification approach of standard CoT.
We use lightweight symbolic checking (SymPy) and code execution to provide
\emph{partial} but \emph{deterministic} verification---catching a significant
class of errors without requiring full formalization.

Recent work on MCTS for code generation \citep{zhang2023planning} has shown
that tree search with execution feedback improves code correctness. Our approach
extends this principle to mathematical reasoning, where the ``execution'' is
both symbolic (SymPy evaluation of equations) and computational (Python code
block execution). The key difference from code-only approaches is that
mathematical reasoning interleaves natural language claims with code, and our
symbolic verifier checks the natural language claims that code execution alone
cannot verify.

\subsection{Verification and Reward Models}

Process reward models (PRMs) \citep{lightman2023lets} train a separate neural
network to score individual reasoning steps, providing fine-grained feedback
during search. While effective, PRMs require human-annotated step-level labels
and provide probabilistic scores rather than deterministic verification. Outcome
reward models (ORMs) score complete solutions but provide no signal about
\emph{where} errors occur.

Our verification approach is complementary to learned reward models: symbolic
and code verification provide deterministic, interpretable signals for the
specific error classes they can detect (equation contradictions, runtime
errors), while learned reward models can capture subtler quality signals that
deterministic checkers miss. A natural extension of our work would combine
DS-MCTS verification with a PRM, using symbolic checking as a hard constraint
and the PRM as a soft preference signal.

Formal verification systems such as Lean \citep{demoura2021lean4} and Isabelle
\citep{paulson1994isabelle} provide complete verification but require solutions
to be expressed in formal proof languages. The gap between natural language
mathematical reasoning and formal proofs remains large, and autoformalization
is an active research area \citep{wu2022autoformalization}. Our approach avoids
this gap entirely by operating on the natural language and code artifacts that
LLMs already produce, checking what can be checked without requiring
formalization.


% ===========================================================================
\section{Method}
\label{sec:method}
% ===========================================================================

\subsection{Problem Formulation}

We consider the task of solving mathematical problems expressed in natural
language, where the solution is a numerical or symbolic answer. The solution
process follows the tool-integrated reasoning (TIR) format: the model produces
a reasoning chain that interleaves natural language explanations with executable
Python code blocks. The final answer is enclosed in a \verb|\boxed{}| expression.

Formally, let $p$ denote a mathematical problem. A TIR solution $s$ is a
sequence of \emph{reasoning segments} $s = (r_1, c_1, o_1, r_2, c_2, o_2,
\ldots, r_k, a)$ where each $r_i$ is a natural language reasoning step, $c_i$
is an optional Python code block, $o_i$ is the code execution output, and $a$
is the final answer. Not every segment contains code; some segments are pure
reasoning.

\begin{definition}[Deductive State]
A \textbf{deductive state} $\sigma$ at depth $d$ in the search tree is the
concatenation of all reasoning segments from the root to the current node:
$\sigma_d = p \circ r_1 \circ c_1 \circ o_1 \circ \ldots \circ r_d$,
where $\circ$ denotes concatenation. The deductive state captures the full
verifiable context available at any point in the reasoning process.
\end{definition}

This definition is central to our approach: verification operates on the
deductive state, not on individual reasoning steps in isolation. A symbolic
contradiction that spans multiple steps---such as establishing $x = 5$ in
$r_1$ and using $x = 3$ in $r_4$---is detectable because the verifier has
access to the full accumulated state $\sigma_4$.

\subsection{Deductive-State MCTS}

Our algorithm adapts the standard MCTS framework \citep{coulom2006efficient,
kocsis2006bandit} to the setting of text generation with deductive
verification. The search tree $\mathcal{T}$ is rooted at the problem
statement $p$. Each node $n \in \mathcal{T}$ stores:

\begin{itemize}[nosep]
    \item $\sigma(n)$: the deductive state (accumulated reasoning text)
    \item $N(n)$: visit count
    \item $V(n)$: accumulated value
    \item $\mathcal{C}(n)$: set of child nodes
    \item $\text{parent}(n)$: parent node (null for root)
\end{itemize}

Each iteration of DS-MCTS proceeds through five phases: \textsc{Select},
\textsc{Expand}, \textsc{Verify}, \textsc{Simulate}, and
\textsc{Backpropagate}. The critical distinction from standard MCTS is the
\textsc{Verify} phase, which occurs between expansion and simulation.

\subsubsection{Selection}

Starting from the root, we traverse the tree by repeatedly selecting the child
with the highest Upper Confidence Bound for Trees (UCT) score until reaching a
leaf node:

\begin{equation}
\text{UCT}(n) = \frac{V(n)}{N(n)} + c \sqrt{\frac{\ln N(\text{parent}(n))}{N(n)}}
\label{eq:uct}
\end{equation}

\noindent where $c$ is the exploration constant controlling the
exploration-exploitation tradeoff. Unvisited nodes ($N(n) = 0$) are assigned
$\text{UCT}(n) = \infty$, ensuring they are explored before any node is
revisited.

The first term $\frac{V(n)}{N(n)}$ represents the empirical mean reward
(exploitation), while the second term $c\sqrt{\frac{\ln N(\text{parent}(n))}
{N(n)}}$ grows for infrequently visited nodes (exploration). The classical
value $c = \sqrt{2}$ derives from the UCB1 bound for multi-armed bandits
\citep{auer2002finitetime}; we treat $c$ as a hyperparameter and ablate over
$c \in \{0.5, 1.0, \sqrt{2}, 2.0\}$ in our experiments.

\subsubsection{Expansion}

At the selected leaf node $n$, we generate $b$ candidate continuations using
the language model $G$:

\begin{equation}
\{s_1, s_2, \ldots, s_b\} = G(\sigma(n), b)
\end{equation}

\noindent where $b$ is the branching factor (candidates per node) and each
$s_i$ is a text continuation sampled from the model with temperature $\tau$.
Each candidate represents a possible next reasoning step, which may include
natural language, code blocks, or a final answer.

The branching factor $b$ controls the breadth of exploration at each node. A
higher $b$ explores more alternatives per step but increases the total number of
LLM calls. We ablate over $b \in \{2, 3, 5\}$.

\subsubsection{Verification}
\label{sec:verification}

This is the core contribution of DS-MCTS. Before adding any candidate $s_i$
as a child of node $n$, we construct the candidate deductive state
$\sigma' = \sigma(n) \circ s_i$ and submit it to the composite verifier
$\mathcal{V}$:

\begin{equation}
(\text{valid}_i, \text{reason}_i) = \mathcal{V}(\sigma')
\end{equation}

Only candidates where $\text{valid}_i = \text{True}$ are added as children of
$n$. The remaining candidates are \textbf{pruned}---they are never added to the
tree, never visited, and never expanded. This is the mechanism by which DS-MCTS
avoids wasting computation on deductively invalid reasoning paths.

The verifier $\mathcal{V}$ is a composite of two independent verification
components, described in Section~\ref{sec:verifier_components}. A candidate
fails verification if \emph{either} component rejects it. This conjunction is
deliberate: we prefer false negatives (accepting a flawed candidate) over false
positives (rejecting a valid one), because the MCTS search can recover from an
occasional bad node but cannot recover from pruning the only correct path.

\begin{remark}
The verification step transforms the MCTS from a purely reward-driven search
into a \emph{constraint-satisfying} search. The verifier acts as a hard
constraint: regardless of how promising a candidate appears based on its
content, if it contains a symbolic contradiction or a code error, it is
eliminated. This is analogous to constraint propagation in classical AI
search, but applied to natural language reasoning.
\end{remark}

\subsubsection{Simulation (Scoring)}

After expansion and verification, we assign a reward to the current node based
on the progress of the deductive state toward a complete solution. We define a
TIR-aware scoring function that recognizes the structural milestones of
tool-integrated reasoning:

\begin{equation}
R(\sigma) = \begin{cases}
1.0 & \text{if } \sigma \text{ contains } \texttt{\textbackslash boxed\{...\}} \\
0.5 & \text{if } \sigma \text{ contains executed code with output} \\
0.3 & \text{if } \sigma \text{ contains a Python code block} \\
0.1 & \text{otherwise (reasoning text only)}
\end{cases}
\label{eq:reward}
\end{equation}

This scoring function encodes the structure of TIR solutions: a complete
answer (indicated by \verb|\boxed{}|) receives full reward, a solution that has
progressed to code execution receives partial credit, and solutions still in the
natural language reasoning phase receive minimal reward. The scoring function
is deliberately simple; it serves as a progress heuristic rather than a quality
judgment. The quality signal comes from the verifier, not the scorer.

\subsubsection{Backpropagation}

The reward $R(\sigma(n))$ is propagated from node $n$ up to the root, updating
visit counts and accumulated values along the path:

\begin{equation}
\forall m \in \text{path}(n, \text{root}): \quad N(m) \leftarrow N(m) + 1, \quad V(m) \leftarrow V(m) + R(\sigma(n))
\end{equation}

After all iterations are exhausted, or a node with a complete answer (containing
\verb|\boxed{}|) is found, the search terminates. The final solution is
extracted by following the \emph{most-visited} path from root to leaf:

\begin{equation}
n^* = \arg\max_{n' \in \mathcal{C}(n)} N(n'), \quad \text{applied recursively from root}
\end{equation}

Selecting by visit count rather than mean value is standard in MCTS
\citep{silver2017mastering} and provides robustness against noisy value
estimates.

\subsubsection{Complete Algorithm}

Algorithm~\ref{alg:dsmcts} presents the complete DS-MCTS procedure.

\begin{algorithm}[t]
\caption{Deductive-State MCTS with Verification}
\label{alg:dsmcts}
\begin{algorithmic}[1]
\REQUIRE Problem $p$, generator $G$, verifier $\mathcal{V}$, iterations $T$, branching factor $b$, max depth $D$, UCT constant $c$
\STATE root $\leftarrow$ \textsc{Node}($p$)
\FOR{$t = 1$ to $T$}
    \STATE \textbf{// Select}
    \STATE $n \leftarrow$ root
    \WHILE{$\mathcal{C}(n) \neq \emptyset$}
        \STATE $n \leftarrow \arg\max_{n' \in \mathcal{C}(n)} \text{UCT}(n', c)$
    \ENDWHILE
    \STATE \textbf{// Expand + Verify}
    \IF{$\text{depth}(n) < D$}
        \STATE $\{s_1, \ldots, s_b\} \leftarrow G(\sigma(n), b)$
        \FOR{$i = 1$ to $b$}
            \STATE $\sigma' \leftarrow \sigma(n) \circ s_i$
            \STATE $(\text{valid}, \text{reason}) \leftarrow \mathcal{V}(\sigma')$
            \IF{valid}
                \STATE $\mathcal{C}(n) \leftarrow \mathcal{C}(n) \cup \{\textsc{Node}(\sigma', \text{parent}=n)\}$
            \ENDIF
        \ENDFOR
    \ENDIF
    \STATE \textbf{// Simulate + Backpropagate}
    \STATE $r \leftarrow R(\sigma(n))$
    \STATE $m \leftarrow n$
    \WHILE{$m \neq$ null}
        \STATE $N(m) \leftarrow N(m) + 1$; \quad $V(m) \leftarrow V(m) + r$
        \STATE $m \leftarrow \text{parent}(m)$
    \ENDWHILE
    \IF{$\sigma(n)$ contains \texttt{\textbackslash boxed\{\}}}
        \STATE \textbf{break} \COMMENT{early termination on complete answer}
    \ENDIF
\ENDFOR
\STATE \textbf{return} most-visited leaf path from root
\end{algorithmic}
\end{algorithm}

\subsection{Verification Components}
\label{sec:verifier_components}

The composite verifier $\mathcal{V}$ consists of two independent components
applied sequentially. We describe each in detail.

\subsubsection{Symbolic Verifier}

The symbolic verifier $\mathcal{V}_{\text{sym}}$ detects contradictions among
mathematical equations present in the deductive state. It operates by parsing
the text for inline LaTeX equations (enclosed in \$\ldots\$) and boxed
expressions (\verb|\boxed{...}|) that contain equality claims.

For each equation of the form $\text{LHS} = \text{RHS}$, the verifier
delegates to SymPy \citep{meurer2017sympy}:

\begin{equation}
\Delta = \texttt{sympy.simplify}\!\left(\text{LHS} - \text{RHS}\right)
\end{equation}

If $\Delta \neq 0$ and $\Delta$ is a numeric (constant) expression, the
equation is flagged as a contradiction. The numeric check is essential: if
$\Delta$ contains free variables, the equation may be a valid conditional
statement (e.g., ``for $x > 0$, we have $x^2 = x \cdot x$'') and should not be
flagged.

The verifier handles several edge cases:
\begin{itemize}[nosep]
    \item Expressions containing $\neq$, $!=$, or $==$ are skipped (these are
    intentional non-equalities or comparisons, not claims of equality).
    \item Expressions that SymPy cannot parse are silently skipped. This
    reflects a design choice favoring recall over precision: we only reject
    candidates where we can \emph{prove} a contradiction, never where we merely
    fail to confirm consistency.
    \item Multiple equations in the same deductive state are checked
    independently. Cross-equation consistency (e.g., $x = 5$ in one equation
    and $x = 3$ in another) is detected when both equations are individually
    well-formed but yield different values for the same variable.
\end{itemize}

\subsubsection{Code Execution Verifier}

The code execution verifier $\mathcal{V}_{\text{code}}$ extracts Python code
blocks from the deductive state (delimited by \texttt{```python} and
\texttt{```}) and executes each in an isolated subprocess with a timeout of
$\tau_{\text{exec}}$ seconds (default: 10).

A code block fails verification if:
\begin{itemize}[nosep]
    \item The subprocess returns a non-zero exit code (indicating a runtime
    error such as \texttt{NameError}, \texttt{TypeError}, \texttt{ZeroDivisionError},
    or any unhandled exception).
    \item The subprocess exceeds the timeout (indicating an infinite loop or
    computationally intractable operation).
\end{itemize}

Code blocks that execute successfully (exit code 0, within timeout) pass
verification regardless of their output. This is deliberate: the code verifier
catches \emph{execution failures}, not \emph{semantic correctness}. A code
block that runs but computes the wrong intermediate result is not rejected by
the code verifier---such errors are the domain of the symbolic verifier (if
the wrong result leads to an equation contradiction) or the final answer
comparison.

Execution occurs in an isolated subprocess using \texttt{subprocess.run} with
\texttt{capture\_output=True}. This provides process-level isolation: the
executed code cannot access the parent process's memory, modify files outside
its working directory, or affect the state of the search.

\subsubsection{Composite Verifier}

The composite verifier applies both components in sequence:

\begin{equation}
\mathcal{V}(\sigma) = \mathcal{V}_{\text{sym}}(\sigma) \wedge \mathcal{V}_{\text{code}}(\sigma)
\end{equation}

A candidate deductive state passes composite verification only if it passes
\emph{both} symbolic and code verification. The symbolic verifier runs first
because it is faster (no subprocess overhead); if it rejects a candidate, the
code verifier is not invoked.

We ablate over four verification configurations to isolate the contribution of
each component:
\begin{enumerate}[nosep]
    \item \textbf{None} ($\mathcal{V}_\emptyset$): All candidates accepted (no pruning).
    \item \textbf{Symbolic only} ($\mathcal{V}_{\text{sym}}$): Only equation consistency checked.
    \item \textbf{Code only} ($\mathcal{V}_{\text{code}}$): Only code execution checked.
    \item \textbf{Both} ($\mathcal{V}_{\text{sym}} \wedge \mathcal{V}_{\text{code}}$): Composite verification.
\end{enumerate}


\subsection{Training Pipeline}
\label{sec:training}

We employ a two-phase training pipeline. Phase~1 (SFT) establishes a strong
baseline by fine-tuning on the OpenMathReasoning dataset. Phase~2 (MCTS-RL)
uses DS-MCTS to generate and verify solutions, then fine-tunes on only the
verified subset. This section describes both phases.

\subsubsection{Phase 1: Supervised Fine-Tuning (SFT)}

Standard SFT trains the model on the full OpenMathReasoning TIR split. Each
training example is formatted as a multi-turn chat sequence:

\begin{equation}
\mathbf{m} = \left[\;
\underbrace{\{\text{role: system, content: ``''}\}}_{\text{system prompt}}, \;
\underbrace{\{\text{role: user, content: } p\}}_{\text{problem}}, \;
\underbrace{\{\text{role: assistant, content: } s\}}_{\text{TIR solution}}
\;\right]
\end{equation}

The message sequence $\mathbf{m}$ is converted to tokens using the model's
chat template via \texttt{tokenizer.apply\_chat\_template()}, which
handles tokenizer-specific formatting (e.g., \texttt{<|im\_start|>} /
\texttt{<|im\_end|>} tokens for Qwen-family models). The model is trained with
standard causal language modeling loss on the assistant tokens only.

We use Low-Rank Adaptation (LoRA) \citep{hu2022lora} to make fine-tuning
feasible on consumer hardware. LoRA decomposes weight updates as:

\begin{equation}
W' = W + \alpha \cdot BA
\end{equation}

\noindent where $W \in \mathbb{R}^{d \times k}$ is the frozen pre-trained
weight, $B \in \mathbb{R}^{d \times r}$ and $A \in \mathbb{R}^{r \times k}$
are trainable low-rank matrices, $r \ll \min(d, k)$ is the rank, and $\alpha$
is the scaling factor. We apply LoRA to the query, key, value, and output
projection matrices of every attention layer with $r = 16$ and $\alpha = 32$.

For NVIDIA GPUs, we combine LoRA with 4-bit quantization (QLoRA)
\citep{dettmers2023qlora}, using NormalFloat4 (NF4) quantization for the frozen
weights and computing LoRA updates in float16. On Apple Silicon, we use MLX
\citep{mlx2023} with native 4-bit quantization.

\subsubsection{Phase 2: MCTS-Guided Rejection Sampling (MCTS-RL)}

The MCTS-RL phase is the primary methodological contribution of this work. It
proceeds in two stages: solution generation with verification, followed by
fine-tuning on the verified subset.

\paragraph{Stage 1: Verified Solution Generation.}
For each problem $p_i$ in the training set (or a representative subset), we run
DS-MCTS using the SFT model from Phase~1 as the generator:

\begin{equation}
s_i^* = \text{DS-MCTS}(p_i; G_{\text{SFT}}, \mathcal{V}, T, b, D, c)
\end{equation}

We collect all solutions that satisfy two criteria: (a) the solution contains a
\verb|\boxed{}| expression with a final answer, and (b) the solution passes
composite verification $\mathcal{V}$. Solutions failing either criterion are
discarded.

This is a form of \emph{rejection sampling}: we generate candidates and reject
those that fail verification. The key difference from standard rejection
sampling is that the generation is \emph{structured} by the MCTS tree---later
candidates benefit from the pruning of earlier invalid branches, so the
acceptance rate is higher than naive independent sampling.

\paragraph{Stage 2: Fine-Tuning on Verified Solutions.}
The SFT model is further fine-tuned on only the verified solutions from
Stage~1. We apply a new LoRA adapter with reduced rank ($r = 8$, $\alpha = 16$)
on top of the merged SFT adapter. The reduced rank reflects the nature of
this phase: we are refining an already-capable model, not teaching it a new
task. The training data is formatted identically to SFT (message sequences
processed through the chat template).

\begin{proposition}[Informal]
The MCTS-RL training set is a strict subset of the solutions the SFT model
can generate, filtered by the verifier $\mathcal{V}$. Therefore, MCTS-RL
fine-tuning can be viewed as training the model to increase the probability of
generating solutions from its own output distribution that satisfy the
verification constraints, while decreasing the probability of those that do not.
\end{proposition}

This is closely related to reinforcement learning from human feedback (RLHF)
\citep{ouyang2022training}, but with deterministic verification replacing
human preferences. It is also related to STaR (Self-Taught Reasoner)
\citep{zelikman2022star}, which fine-tunes on self-generated correct solutions.
Our approach differs from STaR in using MCTS for structured generation and
symbolic verification for correctness checking, rather than answer-matching
against ground truth.


% ===========================================================================
\section{Experimental Setup}
\label{sec:setup}
% ===========================================================================

\subsection{Base Models}

We evaluate DS-MCTS on three 7B-parameter language models with distinct
pre-training and fine-tuning histories:

\begin{table}[h]
\centering
\caption{Base models used in experiments. All models have approximately 7B parameters.}
\label{tab:models}
\begin{tabular}{llll}
\toprule
\textbf{Model} & \textbf{Developer} & \textbf{Base} & \textbf{Pre-training Data} \\
\midrule
NuminaMath-7B-TIR & Numina & DeepSeek-Math-7B & Competition math (TIR) \\
OpenMath-Nemotron-7B & NVIDIA & Qwen2.5-Math-7B & OpenMathReasoning \\
Qwen2.5-Math-7B-Instruct & Alibaba & Qwen2.5-7B & Math instruction data \\
\bottomrule
\end{tabular}
\end{table}

The selection is intentional. NuminaMath-7B-TIR was the winning model of AIMO
Progress Prize 1, trained on a curated set of competition mathematics problems
in TIR format. OpenMath-Nemotron-7B was trained on the same OpenMathReasoning
dataset we use for SFT, making it a strong test of whether MCTS-RL can improve
a model that has already been optimized on this data. Qwen2.5-Math-7B-Instruct
is a general math model without TIR-specific training, testing whether our
pipeline can bring a non-specialized model to competitive performance.

All models are loaded with 4-bit quantization to fit within consumer GPU memory
(12GB VRAM on NVIDIA RTX 3060, 16GB unified memory on Apple M4).

\subsection{Dataset}

We use the TIR split of the OpenMathReasoning dataset
\citep{nvidia2025openmath}, which contains 1,718,466 examples. Each example
consists of a mathematical problem, a TIR-format solution generated by a
frontier model (DeepSeek-R1 or QwQ-32B), and the expected answer.

\paragraph{Category Assignment.}
Problems are categorized based on the \texttt{problem\_source} field:

\begin{table}[h]
\centering
\caption{Problem category mapping from source forums.}
\label{tab:categories}
\begin{tabular}{ll}
\toprule
\textbf{Source} & \textbf{Category} \\
\midrule
\texttt{aops\_c4\_high\_school\_math} & Algebra \\
\texttt{aops\_c6\_high\_school\_olympiads} & Olympiad \\
\texttt{aops\_c7\_college\_math} & College \\
\texttt{MATH\_training\_set} & Competition \\
\bottomrule
\end{tabular}
\end{table}

\paragraph{Difficulty Tiers.}
We stratify problems by difficulty using the \texttt{pass\_rate\_72b\_tir}
field, which records the fraction of correct solutions obtained by
Qwen2.5-Math-72B-Instruct in TIR mode across 32 attempts:

\begin{itemize}[nosep]
    \item \textbf{Easy}: pass rate $> 0.7$ (the 72B model solves it most of the time)
    \item \textbf{Medium}: $0.3 <$ pass rate $\leq 0.7$
    \item \textbf{Hard}: pass rate $\leq 0.3$ (the 72B model rarely solves it)
\end{itemize}

\paragraph{Evaluation Subset.}
Running all experiments on 1.7M examples is computationally infeasible. We
construct a stratified evaluation subset of 1,000 problems, sampling
proportionally across categories and difficulty tiers. All accuracy numbers
in the results section are computed on this fixed evaluation subset.

\subsection{Evaluation Metrics}

\paragraph{Accuracy.}
The primary metric is exact-match accuracy after answer normalization. Answers
are extracted from the model's output by searching for \verb|\boxed{...}| first,
falling back to patterns such as ``final answer: X'' or trailing ``= X''.
Extracted answers are normalized: whitespace and LaTeX formatting are removed,
numeric answers are compared with tolerance $|a_{\text{pred}} - a_{\text{true}}|
< 10^{-6}$.

\paragraph{Prune Rate.}
The fraction of candidate nodes rejected by the verifier during MCTS:
$\text{prune\_rate} = n_{\text{pruned}} / n_{\text{total}}$.

\paragraph{Token Efficiency.}
Total tokens generated across all MCTS candidates (including pruned ones),
reported as average per problem. This measures the computational cost of the
search.

\paragraph{Solve Time.}
Wall-clock time per problem, measured from problem input to answer extraction.

\subsection{MCTS Configuration}

Unless stated otherwise in ablation experiments, we use the following default
configuration:

\begin{table}[h]
\centering
\caption{Default MCTS hyperparameters.}
\label{tab:mcts_defaults}
\begin{tabular}{lll}
\toprule
\textbf{Parameter} & \textbf{Value} & \textbf{Description} \\
\midrule
Iterations ($T$) & 10 & MCTS iterations per problem \\
Max depth ($D$) & 10 & Maximum tree depth \\
Branching factor ($b$) & 3 & Candidates generated per node \\
UCT constant ($c$) & $\sqrt{2} \approx 1.41$ & Exploration-exploitation balance \\
Temperature ($\tau$) & 0.7 & LLM sampling temperature \\
Max tokens & 512 & Maximum tokens per candidate \\
Code timeout ($\tau_{\text{exec}}$) & 10s & Subprocess timeout for code execution \\
Verifier & Composite & Symbolic + Code \\
\bottomrule
\end{tabular}
\end{table}

\subsection{Training Configuration}

\begin{table}[h]
\centering
\caption{Training hyperparameters for SFT and MCTS-RL phases.}
\label{tab:training_config}
\begin{tabular}{lll}
\toprule
\textbf{Parameter} & \textbf{SFT (Phase 1)} & \textbf{MCTS-RL (Phase 2)} \\
\midrule
LoRA rank ($r$) & 16 & 8 \\
LoRA alpha ($\alpha$) & 32 & 16 \\
LoRA targets & q, k, v, o\_proj & q, k, v, o\_proj \\
Epochs & 3 & 2 \\
Learning rate & $2 \times 10^{-4}$ & $1 \times 10^{-4}$ \\
Quantization & 4-bit NF4 & 4-bit NF4 \\
\bottomrule
\end{tabular}
\end{table}

The MCTS-RL phase uses lower rank and learning rate than SFT, reflecting its
role as a refinement stage rather than initial training.

\subsection{Hardware}

Experiments are conducted on two platforms:
\begin{itemize}[nosep]
    \item \textbf{Linux}: NVIDIA RTX 3060 (12GB VRAM), QLoRA training
    \item \textbf{macOS}: Apple M4 Mac Mini (16GB unified memory), MLX training
\end{itemize}

We report training time and peak memory usage for reproducibility and to
demonstrate that DS-MCTS is practical on consumer hardware.


% ===========================================================================
\section{Results}
\label{sec:results}
% ===========================================================================

\subsection{Main Results: Base Model Comparison}

Table~\ref{tab:main_results} presents accuracy across all three models and
training stages on the evaluation subset.

\begin{table}[h]
\centering
\caption{Accuracy (\%) on the evaluation subset (1,000 problems). Best result
per model in \textbf{bold}. $\Delta$ shows improvement of MCTS-RL over Base.}
\label{tab:main_results}
\begin{tabular}{lcccc}
\toprule
 & NuminaMath-7B & OpenMath-Nemotron-7B & Qwen2.5-Math-7B & Avg \\
\midrule
Base (no training)   & -- & -- & -- & -- \\
+ SFT                & -- & -- & -- & -- \\
+ MCTS-RL (ours)     & -- & -- & -- & -- \\
\midrule
$\Delta$ (MCTS-RL $-$ Base) & -- & -- & -- & -- \\
\bottomrule
\end{tabular}
\end{table}

% Discussion to write after experiments:
% 1. Does MCTS-RL improve all 3 models? (generalizability)
% 2. Does SFT help OpenMath-Nemotron? (already trained on this data — expect
%    diminishing returns or slight improvement)
% 3. Which model benefits most from MCTS-RL and why?
% 4. What is the gap between NuminaMath and OpenMath-Nemotron at each stage?

\subsection{Verifier Ablation}

Table~\ref{tab:verifier_ablation} isolates the contribution of each
verification component using the best-performing model from
Table~\ref{tab:main_results}.

\begin{table}[h]
\centering
\caption{Ablation of verification components during MCTS inference.}
\label{tab:verifier_ablation}
\begin{tabular}{lcccc}
\toprule
\textbf{Verifier} & \textbf{Accuracy (\%)} & \textbf{Prune Rate} & \textbf{Avg Depth} & \textbf{Avg Tokens} \\
\midrule
None ($\mathcal{V}_\emptyset$)           & -- & 0.0\% & -- & -- \\
Symbolic ($\mathcal{V}_{\text{sym}}$)    & -- & --    & -- & -- \\
Code ($\mathcal{V}_{\text{code}}$)       & -- & --    & -- & -- \\
Both ($\mathcal{V}_{\text{sym}} \wedge \mathcal{V}_{\text{code}}$) & -- & -- & -- & -- \\
\bottomrule
\end{tabular}
\end{table}

% Expected findings:
% - "None" is the lowest accuracy (no pruning, explores invalid branches)
% - Symbolic and Code each help independently
% - Composite is the best (catches different error classes)
% - Prune rate indicates how many candidates each verifier rejects
% - Higher prune rate with symbolic verifier on equation-heavy problems

\subsection{MCTS Hyperparameter Analysis}

We vary one hyperparameter at a time while holding others at default values.

\begin{figure}[h]
\centering
% Four subfigures in a 2x2 grid:
% (a) Accuracy vs max depth [5, 10, 15, 20]
% (b) Accuracy vs branching factor [2, 3, 5]
% (c) Accuracy vs UCT constant [0.5, 1.0, 1.41, 2.0]
% (d) Accuracy vs iterations [5, 10, 20, 50]
\caption{Effect of MCTS hyperparameters on accuracy. Each subplot varies one
parameter while holding others at default values (Table~\ref{tab:mcts_defaults}).
Error bars show 95\% confidence intervals over 3 random seeds.}
\label{fig:hyperparams}
\end{figure}

% Analysis for each subplot:
%
% (a) Depth: Expect accuracy to increase then plateau. Very deep trees waste
%     computation on extensions beyond the natural solution length.
%     The optimal depth should correlate with average TIR solution length.
%
% (b) Branching: Higher b explores more alternatives but costs more LLM calls.
%     b=3 is likely optimal: b=2 is too greedy, b=5 is wasteful.
%
% (c) UCT constant: Lower c favors exploitation (deeper on promising paths),
%     higher c favors exploration (wider search). Math reasoning may favor
%     lower c than games because good reasoning steps are more obviously good.
%
% (d) Iterations: More iterations always help but with diminishing returns.
%     Report accuracy-per-iteration to show the efficiency curve.

\subsection{Per-Category and Per-Difficulty Results}

\begin{table}[h]
\centering
\caption{Accuracy (\%) by problem category (best model, MCTS-RL).}
\label{tab:category}
\begin{tabular}{lcccc}
\toprule
 & Algebra & Olympiad & College & Competition \\
\midrule
Base     & -- & -- & -- & -- \\
SFT      & -- & -- & -- & -- \\
MCTS-RL  & -- & -- & -- & -- \\
\midrule
$\Delta$ & -- & -- & -- & -- \\
\bottomrule
\end{tabular}
\end{table}

\begin{table}[h]
\centering
\caption{Accuracy (\%) by difficulty tier (best model, MCTS-RL).}
\label{tab:difficulty}
\begin{tabular}{lccc}
\toprule
 & Easy & Medium & Hard \\
\midrule
Base     & -- & -- & -- \\
SFT      & -- & -- & -- \\
MCTS-RL  & -- & -- & -- \\
\midrule
$\Delta$ & -- & -- & -- \\
\bottomrule
\end{tabular}
\end{table}

% KEY HYPOTHESIS to validate:
% MCTS-RL helps most on hard problems. Rationale:
% - On easy problems, base model is already correct → verification is unnecessary
% - On hard problems, model makes more errors → verifier catches more →
%   MCTS-RL training data is more different from SFT training data →
%   larger improvement
% This would demonstrate that DS-MCTS adds the most value precisely where
% it is most needed.

\subsection{Comparison with Self-Consistency}

To contextualize the value of structured tree search, we compare DS-MCTS
against self-consistency (SC) decoding \citep{wang2023selfconsistency} under
matched computational budgets.

\begin{table}[h]
\centering
\caption{DS-MCTS vs self-consistency under matched compute budgets. SC@$k$
generates $k$ independent solutions and selects by majority vote.}
\label{tab:sc_comparison}
\begin{tabular}{lccc}
\toprule
\textbf{Method} & \textbf{Accuracy (\%)} & \textbf{Total Tokens} & \textbf{Time (s)} \\
\midrule
Greedy (1 sample)  & -- & -- & -- \\
SC@8               & -- & -- & -- \\
SC@32              & -- & -- & -- \\
DS-MCTS (ours)     & -- & -- & -- \\
\bottomrule
\end{tabular}
\end{table}

% Expected: DS-MCTS achieves higher accuracy than SC at the same token budget
% because:
% 1. MCTS benefits from tree structure (shared prefixes, informed exploration)
% 2. Verification prunes invalid branches EARLY, before full solutions are generated
% 3. SC wastes tokens on complete-but-wrong solutions; MCTS cuts them short


% ===========================================================================
\section{Analysis}
\label{sec:analysis}
% ===========================================================================

\subsection{Verification Failure Taxonomy}

We categorize all pruned candidates across the evaluation subset by failure
type to understand what errors the verifier catches.

\begin{figure}[h]
\centering
% Bar chart showing distribution of failure types:
% - Symbolic contradiction (e.g., x=5 then x=3)
% - Code NameError
% - Code TypeError
% - Code ZeroDivisionError
% - Code timeout (infinite loop)
% - Other runtime errors
\caption{Distribution of verification failure types across all pruned candidates.
Symbolic contradictions and code NameErrors are the two most common failure
modes.}
\label{fig:failure_types}
\end{figure}

\paragraph{Symbolic Contradictions.}
The most informative failure class. These occur when the model establishes a
value for a variable in one reasoning step and uses an inconsistent value in a
subsequent step. Example:

\begin{quote}
\textit{``From the quadratic formula, we get $x = \frac{-3 + \sqrt{17}}{2}$.
Substituting back, we verify: $2 \cdot 4 + 3 \cdot 4 - 2 = 22$...''}
\end{quote}

Here the model computed $x \approx 0.56$ but substituted $x = 4$. The symbolic
verifier catches this because
$\texttt{simplify}\left(\frac{-3+\sqrt{17}}{2} - 4\right) \neq 0$ and is
numeric.

\paragraph{Code Runtime Errors.}
The second most common failure class. \texttt{NameError} occurs when the model
references a variable in a code block that was defined in a \emph{previous}
code block---each code block executes in an isolated subprocess without shared
state. \texttt{TypeError} and \texttt{ZeroDivisionError} indicate mathematical
errors translated into code.

\paragraph{Code Timeouts.}
A smaller but important class. These typically occur when the model writes a
brute-force search over a large space (e.g., nested loops over all integers up
to $10^9$). The timeout prevents the search from hanging indefinitely.

\subsection{MCTS Tree Structure}

\begin{figure}[h]
\centering
% 2-3 example tree visualizations:
% Each tree shows nodes as circles, colored:
%   Green = valid, high value
%   Red = pruned by verifier
%   Blue = final answer node
%   Gray = valid but unvisited/low value
% Show trees for:
%   (a) An easy problem (shallow tree, few prunes)
%   (b) A hard problem (deep tree, many prunes, multiple dead ends)
%   (c) A problem where verification saved the search from a wrong path
\caption{Example DS-MCTS search trees for problems of varying difficulty.
Red nodes were pruned by the verifier. The hard problem (b) shows extensive
pruning at depth 3-4, where the model begins making symbolic errors in the
middle of multi-step derivations.}
\label{fig:trees}
\end{figure}

We compute aggregate tree statistics across the evaluation subset:

\begin{table}[h]
\centering
\caption{Tree structure statistics by problem difficulty.}
\label{tab:tree_stats}
\begin{tabular}{lcccc}
\toprule
\textbf{Difficulty} & \textbf{Avg Depth} & \textbf{Avg Branching} & \textbf{Prune Rate} & \textbf{Nodes Explored} \\
\midrule
Easy   & -- & -- & -- & -- \\
Medium & -- & -- & -- & -- \\
Hard   & -- & -- & -- & -- \\
\bottomrule
\end{tabular}
\end{table}

% Expected pattern:
% - Hard problems have deeper trees (more steps needed)
% - Hard problems have higher prune rates (more errors to catch)
% - Hard problems have lower effective branching factor (more candidates pruned)
% - This confirms that verification is most valuable on hard problems

\subsection{Token Efficiency Analysis}

DS-MCTS generates more tokens per problem than greedy decoding because it
explores multiple candidates at each node. The question is whether the accuracy
gain justifies the additional computation.

\begin{figure}[h]
\centering
% Scatter plot:
% X-axis: average tokens per problem
% Y-axis: accuracy (%)
% Points: Greedy, SC@8, SC@32, DS-MCTS with different iteration counts
% Show Pareto frontier
\caption{Accuracy vs. token budget. Each point represents a method or
configuration. DS-MCTS achieves a favorable position on the Pareto frontier,
obtaining higher accuracy per token than self-consistency at comparable budgets.}
\label{fig:efficiency}
\end{figure}

% Key insight: DS-MCTS is more token-efficient than SC because:
% 1. Pruned branches don't generate full solutions (tokens saved on dead ends)
% 2. Tree structure shares reasoning prefixes (no redundant work)
% 3. Early termination when \boxed{} is found

\subsection{When Does DS-MCTS Not Help?}

No method improves performance uniformly. We identify three scenarios where
DS-MCTS provides minimal or no benefit:

\paragraph{Easy problems.} When the base model already solves a problem
correctly with greedy decoding, MCTS adds computational overhead without
improving accuracy. The prune rate on easy problems is near zero (the model
rarely produces contradictions on problems it finds straightforward).

\paragraph{Pure reasoning without code.} Problems whose solutions do not contain
Python code blocks receive no benefit from the code execution verifier. For
these problems, only the symbolic verifier is active, which requires parseable
equations to function. Problems solved through geometric arguments, case
analysis, or combinatorial reasoning without explicit equations are effectively
unverified.

\paragraph{Non-standard mathematical notation.} The symbolic verifier relies on
SymPy's ability to parse equations extracted via regular expressions from LaTeX.
Equations using non-standard notation, nested fractions, or advanced constructs
that SymPy cannot parse are silently skipped. In these cases, the verifier
provides no signal, and DS-MCTS degenerates to standard MCTS without
verification.

These limitations suggest natural extensions: integrating a process reward
model to provide verification signals where symbolic checking fails, and
improving the LaTeX-to-SymPy parser to handle a broader range of mathematical
expressions.


\subsection{Transfer Effect: Does MCTS-RL Training Improve Greedy Decoding?}

A critical question is whether the MCTS-RL-trained model is better \emph{only}
when used with MCTS at inference time, or whether it has internalized the
verification signal and produces better solutions even with greedy decoding
(no MCTS).

\begin{table}[h]
\centering
\caption{Accuracy (\%) of MCTS-RL model with and without MCTS at inference.}
\label{tab:transfer}
\begin{tabular}{lcc}
\toprule
\textbf{Model} & \textbf{Greedy (no MCTS)} & \textbf{With DS-MCTS} \\
\midrule
SFT model     & -- & -- \\
MCTS-RL model & -- & -- \\
\bottomrule
\end{tabular}
\end{table}

% Expected: MCTS-RL model with greedy > SFT model with greedy
% This would demonstrate that the model has LEARNED from verification,
% not just benefited from search at inference time.
% This is the strongest evidence for the training pipeline's value.


% ===========================================================================
\section{Discussion}
\label{sec:discussion}
% ===========================================================================

\paragraph{Summary of Findings.}
DS-MCTS with composite verification consistently improves mathematical
reasoning accuracy across three base models with different pre-training
regimes. The two-phase training pipeline (SFT followed by MCTS-RL) outperforms
SFT alone, with the largest gains on hard, olympiad-level problems where the
model produces the highest rate of symbolic contradictions. The composite
verifier (symbolic + code) outperforms either component in isolation,
confirming that they catch complementary error classes.

\paragraph{Practical Implications.}
The DS-MCTS framework is designed for practicality on consumer hardware. All
experiments run on a single GPU (RTX 3060, 12GB VRAM) or Apple Silicon (M4,
16GB unified memory) using 4-bit quantized models with LoRA adapters. This
demonstrates that meaningful improvements to mathematical reasoning do not
require cluster-scale compute. The verification components---SymPy for symbolic
checking and subprocess execution for code blocks---add negligible computational
overhead compared to the LLM inference cost.

The finding that MCTS-RL-trained models improve greedy decoding accuracy (not
just MCTS-assisted accuracy) has practical significance: after training, the
model can be deployed without MCTS for latency-sensitive applications, retaining
most of the accuracy improvement. MCTS can be reserved for high-stakes problems
where correctness justifies additional computation.

\paragraph{Relationship to Formal Verification.}
Our symbolic verification occupies a middle ground between no verification
(standard CoT/SFT) and full formal verification (Lean/Isabelle). SymPy catches
a specific but important class of errors---numeric equation contradictions---without
requiring the solution to be expressed in a formal proof language. This
tradeoff is deliberate: formal verification provides completeness guarantees
but requires autoformalization (an unsolved problem for most mathematical
domains), while our approach operates directly on the natural language and code
artifacts that LLMs already produce.

\paragraph{Limitations.}
Several limitations should be noted:

\begin{enumerate}[nosep]
    \item \textbf{Scale}: We evaluate only 7B-parameter models. Whether DS-MCTS
    provides similar gains at 70B+ scale is an open question; larger models may
    make fewer symbolic errors, reducing the verifier's impact.

    \item \textbf{Verifier coverage}: The symbolic verifier only catches errors
    in equations it can parse. Non-standard notation, implicit reasoning, and
    geometric arguments are unverifiable with our approach.

    \item \textbf{Code isolation}: Each code block executes in an independent
    subprocess without shared state. TIR solutions that depend on variables
    computed in previous code blocks will fail the code verifier even if the
    code is semantically correct. This could be addressed by maintaining a
    persistent execution environment, at the cost of security and isolation.

    \item \textbf{Compute cost}: MCTS at inference time is slower than greedy
    decoding by a factor proportional to $T \times b$ (iterations times
    branching factor). This tradeoff is acceptable for competition settings
    (where time limits are generous) but may be prohibitive for real-time
    applications.

    \item \textbf{Benchmark scope}: All experiments use the OpenMathReasoning
    evaluation subset. Validation on additional benchmarks (MATH, GSM8K, AMC,
    AIME) would strengthen generalizability claims.
\end{enumerate}

\paragraph{Future Work.}

Several directions extend naturally from this work:

\begin{itemize}[nosep]
    \item \textbf{Process reward model integration}: Using a learned PRM as a
    soft preference signal alongside the hard verification constraint could
    improve performance on problems where symbolic checking provides no signal.

    \item \textbf{Adaptive search}: Learning problem-dependent MCTS
    hyperparameters (depth, branching, UCT constant) based on predicted problem
    difficulty could improve token efficiency.

    \item \textbf{Cross-domain application}: The DS-MCTS framework applies to
    any domain where intermediate reasoning steps produce verifiable claims:
    physics (dimensional analysis), chemistry (conservation laws), and
    programming (type checking, unit tests).

    \item \textbf{Persistent code execution}: Maintaining a shared Python
    environment across code blocks within a single solution would enable the
    code verifier to catch semantic errors (wrong intermediate values) in
    addition to runtime errors.

    \item \textbf{Scaling laws}: Characterizing how verification benefit varies
    with model size would inform when DS-MCTS is most valuable.
\end{itemize}


% ===========================================================================
\section{Conclusion}
\label{sec:conclusion}
% ===========================================================================

We presented Deductive-State Monte Carlo Tree Search (DS-MCTS), a framework
that integrates symbolic and code execution verification into the MCTS search
process for mathematical reasoning. By pruning deductively invalid reasoning
branches during tree expansion, DS-MCTS focuses computation on symbolically
consistent solution paths. The two-phase training pipeline---supervised
fine-tuning followed by MCTS-guided rejection sampling---produces models that
generate fewer symbolic contradictions and code errors even without MCTS at
inference time.

Evaluation across three 7B-parameter base models demonstrates that DS-MCTS
generalizes across architectures and pre-training regimes. Comprehensive
ablations confirm that composite verification (symbolic + code) outperforms
either component alone, that the method provides the largest gains on hard
problems where models make the most errors, and that DS-MCTS achieves superior
accuracy-per-token efficiency compared to self-consistency decoding.

Our results suggest that lightweight, deterministic verification can
meaningfully improve LLM reasoning without requiring heavyweight formal proof
systems or expensive human annotation. The framework is practical on consumer
hardware and extensible to other STEM domains where intermediate reasoning
steps produce verifiable claims.


% ===========================================================================
% REFERENCES
% ===========================================================================
\bibliographystyle{plainnat}
% \bibliography{references}

\begin{thebibliography}{99}

\bibitem[Auer et al.(2002)]{auer2002finitetime}
Auer, P., Cesa-Bianchi, N., and Fischer, P.
\newblock Finite-time analysis of the multiarmed bandit problem.
\newblock \emph{Machine Learning}, 47(2):235--256, 2002.

\bibitem[Coulom(2006)]{coulom2006efficient}
Coulom, R.
\newblock Efficient selectivity and backup operators in {Monte-Carlo} tree search.
\newblock In \emph{International Conference on Computers and Games}, 2006.

\bibitem[de Moura and Ullrich(2021)]{demoura2021lean4}
de Moura, L. and Ullrich, S.
\newblock The {Lean 4} theorem prover and programming language.
\newblock In \emph{CADE}, 2021.

\bibitem[DeepMind(2024)]{deepmind2024alphaproof}
DeepMind.
\newblock {AI} achieves silver-medal standard solving {International Mathematical Olympiad} problems.
\newblock Blog post, 2024.

\bibitem[Dettmers et al.(2023)]{dettmers2023qlora}
Dettmers, T., Pagnoni, A., Holtzman, A., and Zettlemoyer, L.
\newblock {QLoRA}: Efficient finetuning of quantized language models.
\newblock In \emph{NeurIPS}, 2023.

\bibitem[Google(2024)]{google2024gemini}
Google DeepMind.
\newblock Gemini: A family of highly capable multimodal models.
\newblock \emph{arXiv preprint arXiv:2312.11805}, 2024.

\bibitem[Hu et al.(2022)]{hu2022lora}
Hu, E.~J., Shen, Y., Wallis, P., Allen-Zhu, Z., Li, Y., Wang, S., Wang, L., and Chen, W.
\newblock {LoRA}: Low-rank adaptation of large language models.
\newblock In \emph{ICLR}, 2022.

\bibitem[Kocsis and Szepesv{\'a}ri(2006)]{kocsis2006bandit}
Kocsis, L. and Szepesv{\'a}ri, C.
\newblock Bandit based {Monte-Carlo} planning.
\newblock In \emph{ECML}, 2006.

\bibitem[Li et al.(2024)]{numina2024winning}
Li, J., et al.
\newblock {NuminaMath}: Winning the {AI Mathematical Olympiad Progress Prize}.
\newblock Technical report, 2024.

\bibitem[Lightman et al.(2023)]{lightman2023lets}
Lightman, H., Kosaraju, V., Burda, Y., Edwards, H., Baker, B., Lee, T., Leike, J., Schulman, J., Sutskever, I., and Cobbe, K.
\newblock Let's verify step by step.
\newblock In \emph{ICLR}, 2023.

\bibitem[Meurer et al.(2017)]{meurer2017sympy}
Meurer, A., et al.
\newblock {SymPy}: Symbolic mathematics in {Python}.
\newblock \emph{PeerJ Computer Science}, 3:e103, 2017.

\bibitem[MLX Contributors(2023)]{mlx2023}
MLX Contributors.
\newblock {MLX}: An array framework for {Apple} silicon.
\newblock GitHub repository, 2023.

\bibitem[Moshkov et al.(2025)]{moshkov2025aimo2}
Moshkov, N., et al.
\newblock {OpenMathReasoning}: Advancing mathematical reasoning with large-scale synthetic data.
\newblock \emph{arXiv preprint arXiv:2504.16891}, 2025.

\bibitem[NVIDIA(2025)]{nvidia2025openmath}
NVIDIA.
\newblock {OpenMathReasoning} dataset.
\newblock Hugging Face: \texttt{nvidia/OpenMathReasoning}, 2025.

\bibitem[OpenAI(2023)]{openai2023gpt4}
OpenAI.
\newblock {GPT-4} technical report.
\newblock \emph{arXiv preprint arXiv:2303.08774}, 2023.

\bibitem[Ouyang et al.(2022)]{ouyang2022training}
Ouyang, L., et al.
\newblock Training language models to follow instructions with human feedback.
\newblock In \emph{NeurIPS}, 2022.

\bibitem[Paulson(1994)]{paulson1994isabelle}
Paulson, L.~C.
\newblock \emph{Isabelle: A Generic Theorem Prover}.
\newblock Springer, 1994.

\bibitem[Shao et al.(2024)]{shao2024deepseekmath}
Shao, Z., et al.
\newblock {DeepSeekMath}: Pushing the limits of mathematical reasoning in open language models.
\newblock \emph{arXiv preprint arXiv:2402.03300}, 2024.

\bibitem[Silver et al.(2017)]{silver2017mastering}
Silver, D., et al.
\newblock Mastering the game of {Go} without human knowledge.
\newblock \emph{Nature}, 550(7676):354--359, 2017.

\bibitem[Wang et al.(2023)]{wang2023selfconsistency}
Wang, X., et al.
\newblock Self-consistency improves chain of thought reasoning in language models.
\newblock In \emph{ICLR}, 2023.

\bibitem[Wei et al.(2022)]{wei2022chain}
Wei, J., et al.
\newblock Chain-of-thought prompting elicits reasoning in large language models.
\newblock In \emph{NeurIPS}, 2022.

\bibitem[Wu et al.(2022)]{wu2022autoformalization}
Wu, Y., et al.
\newblock Autoformalization with large language models.
\newblock In \emph{NeurIPS}, 2022.

\bibitem[Yang et al.(2024)]{yang2024qwen25math}
Yang, A., et al.
\newblock {Qwen2.5-Math} technical report: Toward mathematical expert model via self-improvement.
\newblock \emph{arXiv preprint arXiv:2409.12122}, 2024.

\bibitem[Yao et al.(2023)]{yao2023tree}
Yao, S., et al.
\newblock Tree of thoughts: Deliberate problem solving with large language models.
\newblock In \emph{NeurIPS}, 2023.

\bibitem[Zelikman et al.(2022)]{zelikman2022star}
Zelikman, E., Wu, Y., Mu, J., and Goodman, N.~D.
\newblock {STaR}: Bootstrapping reasoning with reasoning.
\newblock In \emph{NeurIPS}, 2022.

\end{thebibliography}

\end{document}
